\section{Lecture 12 (October 16th)}
\begin{rmk}
Last time we learned about extended isomorphisms. Today, we will look at some specific examples and separable extensions. 
\end{rmk}
\vspace{2ex}
\begin{prop}
(14.1) Let $p(t)\in K[t]$ be a irreducible polynomial and $F/K$ be an extension of fields. $\{$zeros of $p(t)$ in $F$$\}$ is isomorphic to $\{$$K$-embeddings $K[t]/(p(t))\rightarrow F$$\}$.
\end{prop}
\vspace{2ex}

\begin{thm}
(Isomorphism extension theorem) Let $\alpha $ be a zero of an irreducible polynomial $p(t)\in K_1[t]$ and $\beta $ be a zero of $p_{\sigma }(t)\in K_2[t]$. 
\begin{center}
\begin{tikzpicture}[baseline=(current bounding box.center),
  >={Straight Barb[length=3pt,width=6pt]}, % barbed arrowhead (side lines)
  node font=\small]

  % corners
  \node (E) at (0,4.5) {$F_1$};
  \node (F) at (3,4.5) {$F_2$};
  \node (A) at (0,3) {$K_1(\alpha )$};
  \node (B) at (0,0) {$K_1$};
  \node (C) at (3,0) {$K_2$};
  \node (D) at (3,3) {$K_2(\beta )$};

  % vertical arrows (pointing up)
  \draw[->] (B) -- (A);
  \draw[->] (C) -- (D);
  \draw[->] (A) -- (E);
  \draw[->] (D) -- (F);
  \draw[->] (A) -- node[midway,above] {$\exists !\cong$}(D);

  % horizontal arrows (pointing right) with labels above
  \draw[->] (B) -- node[midway, above] {$ \cong$} (C);

  % commutator symbol in the middle (clockwise)

\end{tikzpicture}
\end{center}
One of the concequences of the theorem is the uniqueness of splitting fields: for any non-constant $f(t)\in K[t]$ and splitting fields $F_1$ and $F_2$ for $f(t)$, there exists a $K$-isomorphism
\begin{center}
\begin{tikzpicture}[baseline=(current bounding box.center),
  >={Straight Barb[length=3pt,width=6pt]}, % barbed arrowhead (side lines)
  node font=\small]

  % corners
  \node (A) at (0,0) {$K$};
  \node (B) at (-2,2) {$F_1$};
  \node (C) at (2,2) {$F_2$};

  % vertical arrows (pointing up)
  \draw[->] (B) -- (A);
  \draw[->] (B) -- node[midway,above] {$\cong$}(C);
  \draw[->] (A) -- (C);
\end{tikzpicture}
\end{center}
\end{thm}
\vspace{2ex}
\begin{ex}
(14.13) Let $p(t)=1+t+t^2+t^3+t^{4}\in {\bm Q}[t]$. 
\begin{itemize}
\item[(i)] The splitting field for $p(t)$ is ${\bm Q}(\xi )$
\item[(ii)] $[E:{\bm Q}]=4$
\item[(iii)] The number of ${\bm Q}$-embeddings 
\[E\rightarrow \bar{{\bm Q}}\]
is 4
\end{itemize}
\end{ex}
\vspace{2ex}
\begin{proof}
The (i), note that $(t-1)p(t)=t^{5}-1$. Fix a 5th primitive root of unity (that is, a generator of the multiplicative group $\{z\in {\bm C}\;|\;z^{5}-1=0\}$ which is cyclic). Set $\xi=\exp (2\pi i/5)$ so that $\xi ,\ldots ,\xi ^{4}$ are roots of $p(t)$ in $\bar{{\bm Q}}$. Therefore,
\[E={\bm Q}(\xi ,\xi ^2,\xi ^3,\xi ^{4})={\bm Q}(\xi )\]
For (ii), $[E:{\bm Q}]=[{\bm Q}(\xi ):{\bm Q}]=\mathrm{deg}(\mathrm{irr}(\xi ,{\bm Q}))=\mathrm{deg}(p(t))=4$ by Eisenstein's criterion, which shows that $p(t)$ is irreducible over ${\bm Q}$.

\bigskip

For (iii), we first observe that $E={\bm Q}(\xi )$ is isomorphic to ${\bm Q}[t]/(p(t))$. Applying proposition 14.1, the number of ${\bm Q}$-embeddings from $E={\bm Q}(\xi )\rightarrow \bar{{\bm Q}}$ equals the number of distinct zeros of $p(t)$ in $\bar{{\bm Q}}$, which is 4. 
\end{proof}
\vspace{2ex}
\begin{ex}
(14.14) Let $p(t)=t^3-2\in {\bm Q}[t]$.
\begin{itemize}
\item[(i)] The splitting field $E$ for $p(t)$ is ${\bm Q}(\alpha ,\xi )$.
\item[(ii)] $[E:{\bm Q}]$ is 3.
\item[(iii)] The number of ${\bm Q}$-embeddings 
\[E\rightarrow \bar{{\bm Q}}\]
is 3.
\end{itemize}
\end{ex}
\vspace{2ex}
\begin{proof}
For (i), fix a 3rd primitive root of unity $\xi $. For example, $\xi =\exp (2\pi i/3)$. Then the zeros of $t^3-2$ in ${\bm C}$ (and in $\bar{{\bm Q}}$) are $\alpha ,\alpha \xi $, and $\alpha \xi ^2$, where $\alpha =2^{1/3}$. This implies that 
\[E={\bm Q}(\alpha ,\alpha \xi ,\alpha \xi ^2)={\bm Q}(\alpha ,\xi )\]
For (ii), $[E,{\bm Q}]=[{\bm Q}(\xi ,\alpha ):{\bm Q}]=$ as
\begin{center}
\begin{tikzpicture}[baseline=(current bounding box.center),
  >={Straight Barb[length=3pt,width=6pt]}, % barbed arrowhead (side lines)
  node font=\small]

  % corners
  \node (A) at (0,0) {${\bm Q}$};
  \node (B) at (0,1.5) {${\bm Q}(\xi )$};
  \node (C) at (0,3) {$\big({\bm Q}(\xi )\big)(\alpha )$};

  % vertical arrows (pointing up)
  \draw (A) -- (B);
  \draw (B) -- (C);
\end{tikzpicture}
\end{center}
and
\[\begin{align*}
[{\bm Q}(\xi ,\alpha ),{\bm Q}]=&[{\bm Q}(\xi ,\alpha ):{\bm Q}(\xi )][{\bm Q}(\xi ):{\bm Q}]\\
=&\mathrm{deg}(\mathrm{irr}(\alpha ,{\bm Q}(\xi )))\cdot \mathrm{deg}(\mathrm{irr}(\xi ,{\bm Q}))\\
=&\mathrm{deg}(t^3-2)\cdot \mathrm{deg}(t^2+t+1)=3\cdot 2=6
\end{align*}\]
For (iii), 
\begin{center}
\begin{tikzpicture}[baseline=(current bounding box.center),
  >={Straight Barb[length=3pt,width=6pt]}, % barbed arrowhead (side lines)
  node font=\small]

  % corners
  \node (A) at (0,0) {${\bm Q}$};
  \node (B) at (0,3) {${\bm Q}(\alpha ,\xi )$};
  \node (C) at (3,3) {${\bm Q}(\xi ^2)$};
  \node (D) at (3,0) {${\bm Q}$};
  \node (E) at (3,4.5) {$\bar{{\bm Q}}$};
  \node (F) at (0,1.5) {${\bm Q}(\xi )$};
  \node (G) at (3,1.5) {($\ast$)};

  % vertical arrows (pointing up)
  \draw[->] (F) --node[midway, above] {$2$} (G);
  \draw (A) -- (F); 
  \draw (F) -- (B);
  \draw[->] (B) -- node[midway, above] {$3$} (C);
  \draw (D) -- (G);
  \draw (G) -- (C);
  \draw (C) -- (E);
  \draw[->] (A) -- (D);
  \end{tikzpicture}
\end{center}
Observation is that if $E/K$ and $F/K$ are field extensions, $\sigma :E\rightarrow F$, a $K$-embedding $\alpha \in E$, algebraic over $K$ implies that $\sigma (\alpha )$ is a zero of $\mathrm{irr}(\alpha ,K)$. We make a certain remark. How many ${\bm Q}$-automorhpisms on ${\bm Q}(\xi )$ are there?
\end{proof}
\vspace{2ex}

